% GENERATED FILE. Edit canonical lessons, then run npm run generate:books.
% canonical-source-hash: fnv1a64:605f8981fc2fd5fb
% canonical-lessons: ES-C499-bocadillo, ES-C499-batido, ES-C499-anadir, ES-C499-apetecer, ES-C499-compartir

\chapter{Algo de comer}
\label{ch:es-algo-de-comer}
\hlchaptermodality{499}

\begin{chapteropening}
\textbf{By the end of this chapter:} \emph{I can order a snack, say what I feel like, add something and share.}

Complete the last lesson of chapter 499: I can order a snack, say what I feel like, add something and share.
\end{chapteropening}

\section[el bocadillo]{el bocadillo — the sandwich (in a bread roll)}

\label{lesson:ES-C499-bocadillo}



\begin{quote}
Before the new one: say the Spanish for to ride, then the Spanish for to draw.
\end{quote}

\subsection*{You'll want to know: el bocadillo}
\textbf{el bocadillo} — \textquotedblleft{}the sandwich (in a bread roll)\textquotedblright{}.

\begin{quote}
Un bocadillo de queso, por favor.
\end{quote}

\begin{quote}
\emph{A cheese sandwich, please.}
\end{quote}

The word used in Spain; in much of Latin America a sandwich is \emph{un sándwich}.

\begin{cousinweb}
\textbf{bocadillo} is \emph{bocado}, \textquotedblleft{}a mouthful\textquotedblright{}, with the small-ending \emph{-illo}: a little mouthful. \emph{Bocado} is built on \emph{boca}, \textquotedblleft{}mouth\textquotedblright{}, Latin \emph{bucca}, \textquotedblleft{}a cheek\textquotedblright{} — inherited: the short \emph{u} became \emph{o}.
\end{cousinweb}

\begin{grammarlens}[title={What you've built}]
One more word for reading notices, forms and messages.
\end{grammarlens}

\subsection*{Your turn}
\begin{itemize}
\raggedright
  \item \emph{Say it:} \emph{el bocadillo}
  \item \emph{Say it:} \emph{el bocadillo}, once more
  \item \emph{Say it:} a sentence of your own with \emph{el bocadillo}
  \item \emph{Recall:} say the Spanish for to ride, then the Spanish for to draw, then say \emph{el bocadillo} again
\end{itemize}

\begin{tcolorbox}[breakable,colback=teal!4,colframe=teal!35!black,title={Before you move on}]
What does el bocadillo mean? (\textquotedblleft{}The sandwich (in a bread roll)\textquotedblright{}.) Say it once more.
\end{tcolorbox}
\section[el batido]{el batido — the milkshake, the smoothie}

\label{lesson:ES-C499-batido}



\begin{quote}
Before the new one: say the Spanish for to draw, then the Spanish for the sandwich.
\end{quote}

\subsection*{You'll want to know: el batido}
\textbf{el batido} — \textquotedblleft{}the milkshake, the smoothie\textquotedblright{}.

\begin{quote}
Quiero un batido de fresa.
\end{quote}

\begin{quote}
\emph{I would like a strawberry milkshake.}
\end{quote}

In parts of Latin America it is \emph{el licuado}.

\begin{cousinweb}
\textbf{batido} is the participle of \emph{batir}, \textquotedblleft{}to beat, to whisk\textquotedblright{}, from Latin \emph{battuere}, \textquotedblleft{}to beat\textquotedblright{} (Lewis \& Short): a drink that has been beaten.
\end{cousinweb}

\begin{grammarlens}[title={What you've built}]
One more word for reading notices, forms and messages.
\end{grammarlens}

\subsection*{Your turn}
\begin{itemize}
\raggedright
  \item \emph{Say it:} \emph{el batido}
  \item \emph{Say it:} \emph{el batido}, once more
  \item \emph{Say it:} a sentence of your own with \emph{el batido}
  \item \emph{Recall:} say the Spanish for to draw, then the Spanish for the sandwich, then say \emph{el batido} again
\end{itemize}

\begin{tcolorbox}[breakable,colback=teal!4,colframe=teal!35!black,title={Before you move on}]
What does el batido mean? (\textquotedblleft{}The milkshake, the smoothie\textquotedblright{}.) Say it once more.
\end{tcolorbox}
\section[añadir]{añadir — to add}

\label{lesson:ES-C499-anadir}



\begin{quote}
Before the new one: say the Spanish for the sandwich, then the Spanish for the milkshake.
\end{quote}

\subsection*{You'll want to know: añadir}
\textbf{añadir} — \textquotedblleft{}to add\textquotedblright{}.

\begin{quote}
Añade un poco de sal.
\end{quote}

\begin{quote}
\emph{Add a little salt.}
\end{quote}

\begin{cousinweb}
\textbf{añadir} is usually traced to spoken Latin \emph{*inaddere}, \textquotedblleft{}to add on\textquotedblright{}, built on \emph{addere}, \textquotedblleft{}to add\textquotedblright{} (Lewis \& Short).
\end{cousinweb}

\begin{grammarlens}[title={What you've built}]
One more word for reading notices, forms and messages.
\end{grammarlens}

\subsection*{Your turn}
\begin{itemize}
\raggedright
  \item \emph{Say it:} \emph{añadir}
  \item \emph{Say it:} \emph{añadir}, once more
  \item \emph{Say it:} a sentence of your own with \emph{añadir}
  \item \emph{Recall:} say the Spanish for the sandwich, then the Spanish for the milkshake, then say \emph{añadir} again
\end{itemize}

\begin{tcolorbox}[breakable,colback=teal!4,colframe=teal!35!black,title={Before you move on}]
What does añadir mean? (\textquotedblleft{}To add\textquotedblright{}.) Say it once more.
\end{tcolorbox}
\section[apetecer]{apetecer — to feel like, to fancy}

\label{lesson:ES-C499-apetecer}



\begin{quote}
Before the new one: say the Spanish for the milkshake, then the Spanish for to add.
\end{quote}

\subsection*{You'll want to know: apetecer}
\textbf{apetecer} — \textquotedblleft{}to feel like, to fancy\textquotedblright{}.

\begin{quote}
¿Te apetece un café?
\end{quote}

\begin{quote}
\emph{Do you feel like a coffee?}
\end{quote}

It works like \emph{gustar}: \emph{me apetece}. Used mostly in Spain; elsewhere \emph{tener ganas de}.

\begin{cousinweb}
\textbf{apetecer} is Latin \emph{appetere}, \textquotedblleft{}to reach for, to desire\textquotedblright{} (Lewis \& Short), with the ending \emph{-ecer}. English \emph{appetite} comes from the same verb. The \emph{-t-} between vowels was kept, the mark of a word taken from books.
\end{cousinweb}

\begin{grammarlens}[title={What you've built}]
One more word for reading notices, forms and messages.
\end{grammarlens}

\subsection*{Your turn}
\begin{itemize}
\raggedright
  \item \emph{Say it:} \emph{apetecer}
  \item \emph{Say it:} \emph{apetecer}, once more
  \item \emph{Say it:} a sentence of your own with \emph{apetecer}
  \item \emph{Recall:} say the Spanish for the milkshake, then the Spanish for to add, then say \emph{apetecer} again
\end{itemize}

\begin{tcolorbox}[breakable,colback=teal!4,colframe=teal!35!black,title={Before you move on}]
What does apetecer mean? (\textquotedblleft{}To feel like, to fancy\textquotedblright{}.) Say it once more.
\end{tcolorbox}
\section[compartir]{compartir — to share}

\label{lesson:ES-C499-compartir}



\begin{quote}
Before the new one: say the Spanish for to add, then the Spanish for to feel like.
\end{quote}

\subsection*{You'll want to know: compartir}
\textbf{compartir} — \textquotedblleft{}to share\textquotedblright{}.

\begin{quote}
Compartimos una pizza.
\end{quote}

\begin{quote}
\emph{We share a pizza.}
\end{quote}

\begin{cousinweb}
\textbf{compartir} $\leftarrow$ Late Latin \emph{compartiri}, \textquotedblleft{}to share out\textquotedblright{} (EtymDB 2.1): \emph{con-} and \emph{partiri}, \textquotedblleft{}to divide\textquotedblright{}, built on \emph{pars}, \textquotedblleft{}a part\textquotedblright{}.
\end{cousinweb}

\begin{grammarlens}[title={What you've built}]
One more word for reading notices, forms and messages.
\end{grammarlens}

\subsection*{Your turn}
\begin{itemize}
\raggedright
  \item \emph{Say it:} \emph{compartir}
  \item \emph{Say it:} \emph{compartir}, once more
  \item \emph{Say it:} a sentence of your own with \emph{compartir}
  \item \emph{Recall:} say the Spanish for to add, then the Spanish for to feel like, then say \emph{compartir} again
\end{itemize}

\begin{tcolorbox}[breakable,colback=teal!4,colframe=teal!35!black,title={Before you move on}]
What does compartir mean? (\textquotedblleft{}To share\textquotedblright{}.) Say it once more.
\end{tcolorbox}
